% Stockholm Lerheden CV, prepared 26 October 2026, pending author review.
% OpenAI ChatGPT/Codex: source selection, English profile, ordering and LaTeX.
% Source: data/cv.yaml, data/publications.bib and the shared academic CV.
% Complete publication inventory preserved; no private referee data included.
\documentclass[11pt,a4paper]{article}
\usepackage[a4paper,margin=2.2cm]{geometry}
\usepackage[T1]{fontenc}
\usepackage[utf8]{inputenc}
\usepackage[english]{babel}
\usepackage{lmodern,microtype,amsmath,amssymb,xcolor,enumitem,fancyhdr}
\usepackage[hidelinks]{hyperref}
\setlength{\parindent}{0pt}
\setlength{\parskip}{0.35em}
\setlist[itemize]{leftmargin=1.2em,itemsep=0.15em,topsep=0.2em}
\setlist[enumerate]{leftmargin=1.8em,itemsep=0.5em,topsep=0.2em}
\definecolor{AIReviewRed}{HTML}{B91C1C}
\newcommand{\AI}[1]{{\color{AIReviewRed}#1}}
\pagestyle{fancy}\fancyhf{}\renewcommand{\headrulewidth}{0pt}
\fancyfoot[L]{\small Antonio Michele Miti}
\fancyfoot[R]{\small\thepage}
\begin{document}
\begin{center}
{\LARGE\bfseries Antonio Michele Miti}\\[0.2em]
{\large Academic curriculum vitae}\\[0.3em]
\href{mailto:antoniomichele.miti@gmail.com}{antoniomichele.miti@gmail.com}\quad
\href{https://www.antoniomiti.it/}{antoniomiti.it}\\
\href{https://orcid.org/0000-0002-8829-1943}{ORCID: 0000-0002-8829-1943}\quad 26 October 2026
\end{center}
\section*{Research profile}
\AI{Differential geometer working on multisymplectic observables, higher Dirac and Courant structures, and $L_\infty$-algebras. My joint work develops explicit comparison morphisms and homotopy-fibre models, together with algebraic methods for constraint reduction. The programme connects differential geometry with homotopical algebra and mathematical physics.}
\section*{Research appointments}
\begin{itemize}
\item \textbf{February 2024--April 2026:} CIVIS-3i postdoctoral fellow, Department of Mathematics, Sapienza Universit\`a di Roma. MSCA-H2020 COFUND fellowship on multisymplectic geometry and classical field theories.
\item \textbf{October 2022--January 2024:} Postdoctoral fellow, Universit\`a Cattolica del Sacro Cuore, Brescia. Symplectic and multisymplectic geometry with applications.
\item \textbf{October 2021--September 2022:} Postdoctoral fellow, Max Planck Institute for Mathematics, Bonn. Higher structures and multisymplectic geometry.
\end{itemize}
\section*{Education}
\begin{itemize}
\item \textbf{1 April 2021:} PhD in Mathematics, joint degree of Universit\`a Cattolica del Sacro Cuore and KU Leuven. Thesis: \emph{Homotopy comomentum maps in multisymplectic geometry}. Supervisors: Mauro Spera and Marco Zambon.
\item \textbf{November 2015:} Master's degree in Theoretical Physics, Universit\`a degli Studi di Milano, 110/110. Thesis on algebraic quantization of Jacobi fields and the geometric approach to Peierls brackets. Supervisors: Claudio Dappiaggi and Livio Pizzocchero.
\item \textbf{October 2010:} Bachelor's degree in Physics, Universit\`a degli Studi di Milano-Bicocca. Thesis on Lie groups and rigid body mechanics. Supervisor: Franco Magri.
\end{itemize}
\section*{Selected grants and distinctions}
\begin{itemize}
\item \textbf{2026:} MSCA Seal of Excellence for the GROCOMO-CFT proposal (MSCA-PF-2025).
\item \textbf{2026:} Sapienza \emph{Avvio alla ricerca 2025}, EUR 2,800, supporting the GroCO research project.
\item \textbf{2023:} INdAM research-abroad grant, EUR 3,000.
\end{itemize}
\newpage
\section*{Selected research visits}
\begin{itemize}
\item \textbf{March 2025:} Institut Mittag-Leffler, programme \emph{Cohomological Aspects of Quantum Field Theory}.
\item \textbf{October--December 2024:} Institut Camille Jordan, Universit\'e Claude Bernard Lyon 1, CIVIS-3i secondment.
\item \textbf{December 2023--January 2024:} Institut Camille Jordan, supported by an INdAM research-abroad grant.
\end{itemize}
\section*{Academic service and organisation}
\begin{itemize}
\item Refereeing recorded since 2023 for SIGMA and Geometriae Dedicata.
\item Co-organisation of workshops on multisymplectic geometry: KU Leuven (2020) and Lyon (2024).
\item Assistance with departmental computer systems, Universit\`a Cattolica del Sacro Cuore (2017--2018).
\item \textbf{2024, Lyon:} Organiser, reading seminar \emph{Diffeological Structures in Lagrangian Field Theory}.
\item \textbf{2025, Sapienza:} Co-organiser and host of the doctoral course on perturbative algebraic quantum field theory taught by Fabrizio Zanello.
\end{itemize}
\section*{Teaching and graduate learning activities}
\begin{itemize}
\item \textbf{2024, Salerno:} Co-lecturer with Leonid Ryvkin, research minicourse \emph{Reduction of Multisymplectic Observables}. Eight-hour joint course (total course duration).
\item \textbf{2017--2018, 2022 and 2023, Brescia:} Teaching assistant, \emph{Elements of Higher Geometry}, course led by Mauro Spera. Exercises and selected lectures; support for assessment.
\item \textbf{2018, Brescia:} Teaching assistant, \emph{Differential Topology and Geometry}. Exercises in differential topology and simplicial homology; support for assessment.
\item \textbf{2014--2015, Milano:} Student support at the Physics Department's Computation and Multimedia Laboratory. Computing environments and scientific programming projects, including C, Python and CUDA.
\end{itemize}
\section*{Languages and technical background}
Italian: native. English: TOEFL iBT 102/120, April 2018 (historical test result). French: beginner. Scientific computing and document preparation: Python, C/C++, CUDA, \LaTeX\ and Linux.
\clearpage
\begin{center}{\large\bfseries Publications and research manuscripts}\end{center}
\section*{Published articles and conference papers}
\begin{enumerate}
\item A.~M.~Miti and M.~Zambon. \emph{Observables on multisymplectic manifolds and higher Courant algebroids}. Communications in Contemporary Mathematics \textbf{27} (2025), no.~8, 2550006. Published online in 2024. \href{https://arxiv.org/abs/2209.05836}{arXiv:2209.05836}; \href{https://doi.org/10.1142/S0219199725500063}{DOI}.
\item A.~M.~Miti and L.~Ryvkin. \emph{Multisymplectic observable reduction using constraint triples}. Differential Geometry and its Applications \textbf{100} (2025), 102272. \href{https://arxiv.org/abs/2506.00234}{arXiv:2506.00234}; \href{https://doi.org/10.1016/j.difgeo.2025.102272}{DOI}.
\item C.~Blacker, A.~M.~Miti and L.~Ryvkin. \emph{Reduction of $L_\infty$-algebras of observables on multisymplectic manifolds}. SIGMA \textbf{20} (2024), 061, 47 pp. \href{https://arxiv.org/abs/2206.03137}{arXiv:2206.03137}; \href{https://doi.org/10.3842/SIGMA.2024.061}{DOI}.
\item A.~M.~Miti and L.~Ryvkin. \emph{Multisymplectic actions of compact Lie groups on spheres}. Journal of Symplectic Geometry \textbf{18} (2020), no.~6, 1751--1785. \href{https://arxiv.org/abs/1906.08790}{arXiv:1906.08790}; \href{https://doi.org/10.4310/JSG.2020.v18.n6.a6}{DOI}.
\item A.~M.~Miti and M.~Spera. \emph{A hydrodynamical homotopy co-momentum map and a multisymplectic interpretation of higher order linking numbers}. Journal of the Australian Mathematical Society \textbf{112} (2022), no.~3, 335--354. \href{https://arxiv.org/abs/1805.01696}{arXiv:1805.01696}; \href{https://doi.org/10.1017/S1446788720000518}{DOI}.
\item A.~M.~Miti and M.~Spera. \emph{Derivation of the HOMFLYPT knot polynomial via helicity and geometric quantization}. Bollettino dell'Unione Matematica Italiana \textbf{14} (2021), no.~2, 269--284. \href{https://arxiv.org/abs/1910.13400}{arXiv:1910.13400}; \href{https://doi.org/10.1007/s40574-020-00254-5}{DOI}.
\item A.~Luongo, A.~M.~Miti, V.~Narasimhachar and A.~Sireesh. \emph{Measurement-based uncomputation of quantum circuits for modular arithmetic}. In \emph{2025 62nd ACM/IEEE Design Automation Conference (DAC)}, pp.~1--7. \href{https://doi.org/10.1109/DAC63849.2025.11132981}{DOI}. Expanded preprint: \href{https://arxiv.org/abs/2407.20167}{arXiv:2407.20167}.
\end{enumerate}
\section*{Public preprint}
D.~Fiorenza and A.~M.~Miti. \emph{$L_\infty$-morphisms between twisted Courant $r$-Lie algebras and untwisted Courant $(r+1)$-Lie algebroids}. Preprint (2026), \href{https://arxiv.org/abs/2602.14702}{arXiv:2602.14702}.
\section*{Doctoral thesis}
A.~M.~Miti. \emph{Homotopy comomentum maps in multisymplectic geometry}. Joint doctorate, Universit\`a Cattolica del Sacro Cuore and KU Leuven, awarded 1 April 2021. \href{https://arxiv.org/abs/2105.05645}{arXiv:2105.05645}.
\section*{Manuscripts in preparation}
\begin{itemize}
\item A.~Basu, D.~Fiorenza and A.~M.~Miti. \emph{A morphism between the higher Dirac and higher Courant $L_\infty$-algebras}.
\item M.~V.~D'Agostino, A.~Frabetti and A.~M.~Miti. \emph{Poincar\'e groupoid symmetries and gauge connections}.
\item A.~M.~Miti and F.~Zanello. \emph{Multisymplectic approach to gauge natural Lagrangian field theories and the stress-energy-momentum tensor}.
\end{itemize}

\end{document}
